% ---- pg-516 （书页 504） ----

对任何两个或三个变量的数据在变量空间的分组情况, 都可以用肉眼来识别. 用这种方法往往能识别出潜在的重要分类变量, 并能揭示那些“离群的”或“非典型的”观测值. 可视性显示在判别和分类中是重要的辅助工具, 因而随着计算机硬件和有关计算机程序变得随手可得, 它们的应用也可能不断增加. 从视觉考察所获得的信息往往与通过复杂数值分析所得到的一样多.

\mvsub{关于多元正态性的实际考虑}

尝试性假定的选择与所得到的分离函数形式之间的相互作用是重要的. 考虑\hyperref[fig:11.21]{图 11.21}, 此图显示两个非正态密度的肾形轮廓线. 在这种情况下, 基于正态理论的线性(甚或二次)分类法则与另外一种选择相比就显得不适当了. 也就是说, 线性判别法在此处不适用.

\begin{mvfigure}
\begin{tikzpicture}[scale=0.95,transform shape]
% 坐标轴
\draw[->] (0,0)--(7.3,0) node[below right] {$x_1$};
\draw[->] (0,0)--(0,5.3) node[left] {$x_2$};
% 两个非正态总体的肾形轮廓线
\draw[thick] plot[smooth cycle,tension=0.9] coordinates {
(1.5,1.7) (1.7,2.5) (2.3,3.15) (3.0,3.35) (3.45,3.0)
(3.05,2.7) (2.6,2.45) (3.05,2.15) (3.5,1.9) (3.3,1.4)
(2.7,1.15) (2.0,1.2)};
\draw[thick] plot[smooth cycle,tension=0.9] coordinates {
(3.45,2.95) (4.05,3.6) (4.85,3.6) (5.35,3.1) (5.4,2.3)
(5.05,1.65) (4.35,1.3) (3.7,1.35) (3.35,1.7) (3.85,1.95)
(4.3,2.3) (3.9,2.65)};
% 分界线
\draw[dashed] (1.9,5.05)--(6.4,0.5);
\draw[dashed] (2.75,5.15)--(3.75,0.15);
% 标注
\node at (2.15,4.9) {\small “线性分类”的分界线};
\draw[->] (2.0,4.62)--(1.85,4.35);
\node at (4.6,4.6) {\small “良好分类”的分界线};
\draw[->] (3.95,4.38)--(3.35,4.12);
\node[align=center] at (0.62,2.95) {$f_2(x)$的\\轮廓线};
\draw[->] (1.55,3.0)--(1.85,2.9);
\node at (6.35,3.2) {\small $f_1(x)$的轮廓线};
\draw[->] (5.9,3.1)--(5.4,2.95);
\node at (1.5,0.85) {$R_2$};
\node at (5.95,1.4) {$R_1$};
\end{tikzpicture}
\phantomsection\label{fig:11.21}\mvcap{图 11.21\quad 线性判别不适用的两个非正态总体}
\end{mvfigure}

经常要对有很多变量的对象进行分类, 其中有些是“有—无”型或“0—1”型变量. 在这些情形或其他一些变量值域受限制的情形下, 多元正态性也许不是一个明智的假定. 如我们所知, 基于费希尔线性判别量的分类仅当多元正态性假定成立时才是最小 ECM 或最小 TPM 意义下的最优分类. 当正态性明显不成立时, 我们如何解释这些量?

在没有多元正态性的场合, 费希尔的线性判别量可以看成是对总样本信息提供了一种近似. 前几个判别量本身的正态性可以检验, 然后可以用法则(11-67). 由于判别量是大量变量的线性组合, 它们往往接近正态变量. 当然, 必须记住: 前几个判别量为原始样本信息的不完全汇总. 当根据全部样本信息导出的最优法则表现很好时, 基于这个受限制的判别量集合的分类法则也许表现很差.

\mvsub{练习}

\noindent 11.1\quad 考虑两个数据集
\[
\mathbf{X}_1=\begin{bmatrix}3&7\\2&4\\4&7\end{bmatrix}
\qquad\text{和}\qquad
\mathbf{X}_2=\begin{bmatrix}6&9\\5&7\\4&8\end{bmatrix}
\]
其中
\[
\bar{\mathbf{x}}_1=\begin{bmatrix}3\\6\end{bmatrix},\qquad
\bar{\mathbf{x}}_2=\begin{bmatrix}5\\8\end{bmatrix}
\]
及
